\section{Related work, and what is claimed as new}\label{sec:related}

Three literatures run close to the results of Sections~\ref{sec:general}--\ref{sec:witness}, and the separation is new only in a specific and bounded sense. This section states that sense exactly, so that the reader can tell the transported ingredients from the assembled object.

\subsection{The scalarization gap in multi-objective optimisation}

It is a textbook fact that weighted-sum scalarization recovers only the \emph{supported} non-dominated outcomes -- those on the boundary of the convex hull of the non-dominated set -- and that no choice of weights recovers unsupported points on non-convex portions of the front \citep{Koski1985,AthanPapalambros1996,DasDennis1997,Miettinen1999,Ehrgott2005,MessacMattson2002}. The remedy literature -- adaptive weighted sums, Tchebycheff and $\epsilon$-constraint scalarizations -- exists because of it. Theorem~\ref{thm:menu}(ii) reproduces this geometry exactly: scalarized acceptance is $\conv(\mathcal D)+\Rp$ and the acceptance gap is its excess over the union of the plans' orthants. The convex-hull characterisation should therefore not be read as a new fact about scalarization; it is the scalarization gap, transported from the recovery of Pareto points under a single optimisation to the acceptance of transitions under a quantifier over plans. What the transport changes is the object compared: the multi-objective literature asks whether a weighted-sum optimiser can \emph{reach} an unsupported point, whereas Theorem~\ref{thm:witness} asks whether a transition \emph{certified by every weighting} is implementable at all. The noncompensatory composite-indicator construction of \citet{MundaNardo2009} is the index-side response to the same geometry at a fixed profile; the present results are one level up, in the acceptance protocol.

\subsection{Adjustable and randomised robust optimisation}

Robust optimisation distinguishes the robust counterpart, in which every decision is here-and-now, from the adjustable counterpart, in which a second stage responds to the realised uncertainty \citep{BenTal2004,BenTal2009}. Under constraint-wise uncertainty the two coincide for linear problems; \citet{BertsimasGoyal2010} bound how far a static solution may lie from an adjustable one, and \citet{WeiZhang2024} give a necessary and sufficient condition, in the form of a theorem of the alternatives, for the adjustability gap to vanish on a polyhedral uncertainty set. Their gap is measured in optimal objective \emph{value}; the gap studied here is one of \emph{feasibility} -- whether any common implementable plan exists -- and their polyhedrality hypothesis has no counterpart here. The two bodies of results describe the same quantifier inversion from opposite faces: they characterise when it closes, Theorem~\ref{thm:witness} exhibits when it does not, on a region of nonempty interior. The here-and-now/wait-and-see vocabulary is nevertheless the right name for the protocols of Table~\ref{tab:protocols}: common-plan acceptance is the static robust problem and weight-adaptive acceptance the adjustable one, with the weight playing the role of the revealed uncertainty.

Separately, it is well established that randomising over decisions can strictly improve worst-case performance against a deterministic choice; robust optimisation classifies problems as randomisation-receptive or randomisation-proof according to whether a randomised decision strictly beats the deterministic optimum \citep{Delage2019,KrauseRoperGolovin2011,VorobeychikLi2014,Sinha2018,Sessa2020}. Theorem~\ref{thm:blend} -- admitting fractional blends closes the action-indexed gap exactly -- sits in this line, and the datum is, in that vocabulary, randomisation-receptive under the action-indexed class and randomisation-proof at the canonical point under the shared-shock class (Proposition~\ref{prop:commonshock}(iii)).

A reader from game theory will reach for Kuhn's theorem here, and the point should be met directly. In an extensive-form game with perfect recall, mixed strategies (a distribution over plans drawn once) and behavioural strategies (an independent randomisation at each decision point) are realisation-equivalent: they induce the same distribution over terminal nodes \citep{Kuhn1953}. Theorem~\ref{thm:blend} and Proposition~\ref{prop:timeshare} together assert that blending over plans closes the gap while full-depth alternation between plans over time does not, which on its face contradicts that equivalence. It does not, for a reason that is the paper's own scope restriction. Kuhn's equivalence is an equality of \emph{distributions over paths} evaluated under an expected criterion; the certified constraint here is \emph{path-wise} and must hold on every disturbance trajectory of the class. Under a worst-case criterion the two notions of randomisation are not interchangeable, because convexification must occur \emph{within} a time step -- a blend applied at each instant yields the convex combination of required margins at that instant, whereas full-depth alternation varies \emph{across} steps and yields a vertex at any given step; since the constraint binds at every step, across-step variation at plan granularity never convexifies. Arbitrarily fine duty cycling is a different control class, treated in Section~\ref{sec:closing} through relaxed controls, and is not ruled out by the alternation result. The distinction is reinforced from the other side: \citet{MainRandour2024} show that Kuhn's equivalence fails outright once strategies are restricted to finite memory, which is the setting here.

\subsection{Weak and strong sustainability}

The distinction the typed and compensatory operators formalise is the one environmental economics draws between weak sustainability, under which natural and manufactured capital are substitutes and only the aggregate need be maintained, and strong sustainability, under which critical natural capital is non-substitutable and must be maintained separately \citep{PearceAtkinson1993,Beckerman1994,Daly1995,Ayres2001,Neumayer2013,DietzNeumayer2007}. That debate is conducted largely in verbal and policy terms; it has not, to our knowledge, been given a separation theorem with an exact witness. The maxim advanced in Section~\ref{sec:design} -- that where floors bind separately, reporting them individually is a detection requirement rather than a presentation preference -- is a consequence of the formal result, not a restatement of the policy position, and the ladder of Theorem~\ref{thm:ladder} is the elasticity-side statement of the same debate: limited substitutability shifts the gap but does not remove it.

\subsection{What is claimed as new}\label{sec:whatisnew}

Given the above, the contribution is deliberately narrow. Four things are claimed, and nothing more.
\begin{enumerate}[leftmargin=2em]
\item \textbf{The setting.} The separation is established for \emph{robust transition safety over a finite horizon with path-wise constraints}, not for the recovery of Pareto points by a scalarized optimisation. The objects compared are acceptance \emph{sets of states} under a viability constraint; the argument turns on quantifier order in that setting (Lemmas~\ref{lem:quant}--\ref{lem:cone}, Theorem~\ref{thm:menu}), with the synchronisation hypothesis of Lemma~\ref{lem:sync} fixing which path protocol the margin geometry describes exactly.
\item \textbf{The negative result on time-alternation, which inverts the usual ordering.} Theorem~\ref{thm:blend} and Proposition~\ref{prop:timeshare} show that fractional blending closes the action-indexed gap exactly while full-depth alternation over time does not. In robust optimisation the ordering runs the other way: adjustability weakens conservatism, so the adjustable counterpart dominates the static one. Here convexification in \emph{plan space} achieves what open-loop variation over \emph{time} at plan granularity cannot. We have not found this comparison, or this inversion, in the robust-optimisation or randomisation literatures surveyed above, which study worst-case objective value rather than path-wise feasibility. This is the sharpest claim in the paper and the one most likely to be contested; its scope is the action-indexed class, since under the shared-shock class blending itself fails at the canonical point.
\item \textbf{The detection reading.} Because the gap has nonempty interior on an explicit datum, an aggregate index cannot distinguish the rescuable from the impossible without a separate check of the floors and the reserve, or a verified common implementable plan (Theorem~\ref{thm:witness}, Proposition~\ref{prop:rescue}). This is a statement about what a reported index can and cannot certify, and it is a consequence of items 1 and 2, not an independent assertion.
\item \textbf{The witness and its robustness.} The datum is realised in exact rational arithmetic as a fishery transition under a heatwave disturbance with units anchored to a regulated biomass limit (Section~\ref{sec:benchmark}); the gap is a region of nonempty interior rather than a measure-zero boundary phenomenon; and, on the witness-defined diagonal gap slice, every positive-elasticity CES member falsely certifies an exact interval with algebraic critical floors (Theorem~\ref{thm:ladder}); the gap also survives the replacement of the action-indexed disturbance class by the specified shared shock (Proposition~\ref{prop:commonshock}).
\end{enumerate}
We do \emph{not} claim the quantifier-order observation itself (Lemma~\ref{lem:quant} is elementary order theory), the convex-hull geometry (it is the known scalarization gap), or the usefulness of randomisation under worst-case objectives (well established). A reader familiar with any of the three literatures will recognise the ingredients; the claim is about the specific object they assemble here and, principally, about item~2.

